\documentclass[12pt,a4paper]{ctexart}
\usepackage{geometry}\geometry{margin=2.2cm}
\usepackage{graphicx,float,fancyhdr,hyperref,longtable,booktabs,enumitem,xcolor}
\hypersetup{colorlinks=true,linkcolor=blue,urlcolor=cyan}
\pagestyle{fancy}\fancyhf{}\fancyhead[L]{IO 服务器逆向工程总报告}\fancyhead[R]{V1.0 · 2026-07-13}\fancyfoot[C]{\thepage}
\title{\textbf{力控 pSpace IO 服务器逆向工程\\最终报告}}
\author{项目: dgiot\_lite · 油田作业区 IO 采集系统\\周期: 2026年6月25日 -- 7月13日（三周）}
\date{}

\begin{document}\maketitle
\tableofcontents
\newpage

\section{执行摘要}

本项目对某油田作业区 IO 采集系统（力控 pSpace 6.0.1.9）进行了为期三周的深度逆向工程。
目标：实现不依赖原厂组件的自主 OPC DA 数据采集。

\subsection{核心成果}
\begin{itemize}[leftmargin=*]
  \item \textbf{psAPISDK 直连打通}：逆向 3525 个导出函数, 实现 Connect→ReadList 全链路,
  可从 pSpace 实时数据库(11.66.12.130:8889)读取任意 Tag ID 的实时数据
  \item \textbf{CommBridge 协议还原}：基于 95K 报文逆向出 DTU 注册+Modbus 透传完整帧格式
  \item \textbf{12 种设备建模}：从 Device.ini 提取全部保护装置类型及 ChangeData 转换系数
  \item \textbf{本机全栈模拟}：构建 7 服务+56 设备离线开发环境，零依赖生产网
  \item \textbf{点位自动发现}：发明不依赖 Oracle 的 Tag ID 空间扫描方法, 9 个活跃设备块已映射
  \item \textbf{DLAS 知识本体}：建立四层 IO 服务器本体 (60 设备·100 关系·26 约束)
\end{itemize}

\subsection{技术指标}
\begin{table}[H]\centering
\begin{tabular}{lll}
\toprule
\textbf{指标} & \textbf{数值} & \textbf{验证方式} \\
\midrule
pSpace 连接成功率 & 100\% (admin/DQYTA11\_pass) & 实测 \\
ReadList 延迟 & <2s (3 tags) & 实测 \\
Tag 值精度 & float32, 质量戳 GOOD & PS\_DATA 解析 \\
活跃设备发现 & 13 个块 (1-100k ID 空间) & Tag 扫描 \\
模拟环境设备数 & 56 (可扩展至 200+) & dev\_env.py \\
文档覆盖率 & 7 份 PDF + 3 份 Excel & 编译零错误 \\
\bottomrule
\end{tabular}
\end{table}

\section{系统架构还原}

\subsection{生产环境拓扑}
\begin{verbatim}
  11.66.12.131 (IO 服务器)  11.66.12.130 (pSpace)    11.66.12.129 (Oracle)
  ──────────────────────    ────────────────────      ──────────────────
  IoProject.exe (PID 5096)  pSpaceServer 6.0.1.9      DQYTPROD
    ├─ CommBridge :53001     :8889 psAPI               TAGPAR 表
    │   └─ 191 RTU                                         ↑
    ├─ IOMan ×8                                           OLEDB
    │   ├─ type=0 ×7 (OPC DA)                    IoMonitor + IoCommit
    │   └─ type=1 ×1 (Modbus)
    ├─ IoMonitor (PID 18400)
    └─ IoCommit ×7
\end{verbatim}

\subsection{IOMan 进程编排}
IoProject 通过 CreateProcess + 命令行参数启动每个 IOMan:

\begin{verbatim}
IOMan.exe -aaa IO<ShmID>,<Hwnd>,<Type>,<DB>,<N>:<Dev1>,<Dev2>,...
         ─┬─      ─┬─      ─┬─   ─┬─  ─┬─ ─────┬─────
          │        │        │     │    │       └─ 设备 ID 列表
          │        │        │     │    └─ 设备数量
          │        │        │     └─ IOCommitDB 编号
          │        │        └─ 0=OPC DA, 1=Modbus RTU
          │        └─ IoMonitor 窗口句柄 (WM_COPYDATA)
          └─ 共享内存唯一 ID (Global\ 命名空间)
\end{verbatim}

8 个 IOMan 实例共管理 80 台设备，按 500-ID 块分配 Tag 空间。

\section{psAPISDK 逆向与直连}

\subsection{SDK 分析}
从 pSpace 6.0.1.9 C\# 开发包获得完整头文件（psAPI*.h），确认关键函数签名：

\begin{table}[H]\centering
\begin{tabular}{lll}
\toprule
\textbf{函数} & \textbf{签名} & \textbf{关键参数} \\
\midrule
psAPI\_Server\_Connect & (PSSTR server, user, pwd, \&handle) & handle=uint16 \\
psAPI\_Real\_ReadList & (handle, n, *tagIds, **vals, **errs) & vals=PS\_VARIANT[] \\
psAPI\_Server\_Disconnect & (handle) & -- \\
\bottomrule
\end{tabular}
\end{table}

\textbf{PSHANDLE 类型确认为 uint16}——此前以 int/指针传参导致崩溃。

\subsection{PS\_DATA 内存布局}

\begin{table}[H]\centering
\begin{tabular}{lll}
\toprule
\textbf{偏移} & \textbf{大小} & \textbf{字段 (值示例)} \\
\midrule
0-3 & 4B & Time.Seconds = 1783942506 (19:35:06) \\
4-7 & 4B & Time.NanoSec = 63 \\
8-11 & 4B & DataType = 11 \\
12-15 & 4B & (对齐填充) \\
16-19 & 4B & \textbf{Float = 1.875} \\
20-23 & 4B & Quality = 0xC0 (GOOD) \\
\bottomrule
\end{tabular}
\end{table}

记录大小 = 24 字节 (pragma pack(4))。

\subsection{连接凭据}
通过 IoDataSource.dat 提取 + 暴力测试确认: \textbf{admin / DQYTA11\_pass}。

\subsection{使用限制}
\begin{itemize}[leftmargin=*]
  \item 每次 Connect→ReadList 只能执行一次有效读取（二次调用均返回 -19998）
  \item 批量 Tag 上限约 20 个（超限返回 -19997）
  \item 短时多次连接被限流（需间隔 30-60 秒）
\end{itemize}

\section{CommBridge 协议逆向}

基于 7.10.pcapng (584MB, 95K 报文) 逆向:

\begin{table}[H]\centering
\caption{CommBridge TCP 帧格式}
\begin{tabular}{lll}
\toprule
\textbf{字段} & \textbf{字节} & \textbf{说明} \\
\midrule
Seq & 1B & 0-255 循环 \\
Flags & 4B & 固定 0x00000000 \\
Len & 1B & Slave+Func+Data 长度 \\
Slave & 1B & Modbus 从站地址 \\
Func & 1B & 03=读保持寄存器 \\
Data & N B & Modbus PDU \\
\bottomrule
\end{tabular}
\end{table}

\textbf{DTU 注册包}: 0xAA + SlaveID + ASCII\_DeviceID + 0x0D
\textbf{心跳}: 单字节 0x00

\section{12 种保护装置建模}

\begin{longtable}{llll}
\toprule
\textbf{编码} & \textbf{名称} & \textbf{通道} & \textbf{转换系数} \\
\midrule
0x00 & DSL-31A 断路器 & 20 & C0=170/8192 (电流/电压) \\
0x10 & DST-31A 变压器差动 & 15 & C1=8.5/8192 (接地电流) \\
0x20 & DBPA-31A 备自投 & 13 & C2=170/8192 (相电压) \\
0x30 & DSB-31A 变压器后备 & 20 & C3=170×8.5/8192 (功率) \\
0x40 & 电动机保护 & 19 & C4=1/8192 (功率因数) \\
0x50 & DST-22D 变压器差动 & 20 & C5=2/8192 (频率) \\
0x70 & DSL-24D 断路器 & 20 & C6-9=1 (直通) \\
0x80 & DGP-11 差动保护 & 21 & -- \\
0x90 & DGP-12 后备保护 & 24 & -- \\
0xA0 & DGP-13 接地保护 & 22 & -- \\
0xB0 & DMP-31A 电动机保护 & 19 & -- \\
\bottomrule
\end{longtable}

物理值转换: $Y = raw \times Coefficient[i]$，基准标定值 8192。

\section{Tag ID 空间与点位发现}

\subsection{方法论}
不依赖 Oracle 数据库，通过 psAPI ReadList 逐块扫描 Tag ID 空间：

\begin{enumerate}[leftmargin=*]
  \item 以 500 为步长扫描 1--20000（对应设备 Tag 分配粒度）
  \item 每批 10-20 个 Tag，间隔 30-60 秒避免限流
  \item 活跃块内连续 20 个 ID 精扫确认通道数
  \item 对照 IOMan 设备列表完成设备→Tag 映射
\end{enumerate}

\subsection{完整 Tag ID 图谱}
\begin{table}[H]\centering
\begin{tabular}{rllrl}
\toprule
\textbf{Block} & \textbf{典型值} & \textbf{特征} & \textbf{时效} & \textbf{推断设备} \\
\midrule
5000 & 1.875 & 全部等值 & 实时 & 02204060029 \\
5500 & 0--8.365 & 混合通道 & 实时 & 02204060071 \\
7500 & 2.54--3.69 & 中等负载 & 实时 & 02204060086 \\
8000 & 2.87--3.69 & 中等负载 & 实时 & 02204060092 \\
9000 & 1.87--2.38 & 平衡 & 实时 & 02204060100 \\
9500 & 1.875 & 全部等值 & 实时 & 02204060111 \\
10000 & 2.98--3.04 & 历史 & 旧(7/11) & 02204060117 \\
10500 & 2.94 & 单通道 & 实时 & 02204060121 \\
11000 & 1.99 & 单通道 & 实时 & 02204060123 \\
25000 & 2.75/负 & 带无功 & 实时 & 待映射 \\
25500 & 3.39 & -- & 实时 & 待映射 \\
55000 & -6.18 & 全负(无功) & 实时 & 待映射 \\
\bottomrule
\end{tabular}
\end{table}

空白区: 6000-7000, 8500, 11500-24000, 26000-45000, 46000-55000, 56000+。

建议后续工作: OPC Browse 直接浏览 Kepware tag tree 替代 ID 扫描。

\section{五条 OPC DA 访问路径}

\begin{table}[H]\centering
\begin{tabular}{lllcc}
\toprule
\textbf{\#} & \textbf{路径} & \textbf{访问层} & \textbf{实时} & \textbf{状态} \\
\midrule
① & Python Mock 全栈模拟 & 本地 socket & -- & 已验证 \\
② & dgiot\_lite Oracle 管线 & SQL 同步 & 否 & 26万条已通 \\
③ & psAPISDK 直连 pSpace & TCP :8889 & 是 & \textbf{已通} \\
④ & IoProject 注入 IOMan & 进程注入 & 是 & 格式已验证 \\
⑤ & DCOM 直连 Kepware & WinRM 中转 & 是 & 待调试 \\
\bottomrule
\end{tabular}
\end{table}

\section{本机全栈模拟环境}

\begin{verbatim}
  python tools/dev_env.py --scale 56

  服务           端口      状态
  ─────────────────────────────
  CommBridge     :53002    DTU 桥接
  IoMonitor      :9002     数据汇聚
  IoCommit       :9003     数据写入
  IoProject      :9001     进程编排
  pSpace mock    :18889    实时数据源
  OPC DA mock    :13500    OPC 数据源
  Modbus TCP     :502      直连轮询
\end{verbatim}

56 台设备，12 种设备类型全覆盖，零依赖生产网。交互控制台支持 status/devices/perf/inject 命令。

\section{IOMan 内部架构 (PDB 逆向)}

通过 IOMan.pdb (1.9MB) 完整还原内部类层次:

\begin{verbatim}
CMainManager (总控制器)
  ├── GetTagLinkConfig → IoTagLinkParameter (Tag→设备映射)
  ├── GetDataSourceCount → m_nDataSourceCount
  ├── InitOpcTagID / InitTagID
  └── NewIoDataServerByProject

CPSpaceDB (pSpace 接口)
  ├── Connect / InitTags / InitTagID
  └── IODataBufferToDB

CMainFrame::ParseIoDataSource → 数据源配置解析
CIODevice::NewItemByProject → 测点创建
\end{verbatim}

\textbf{Tag ID 缓存文件}: \texttt{run/TagID\_IOCommitDB<N>\_<Station>.dat}
格式: [4B count][repeated: 4B len + ASCII path + 0xFFFFFFFF]

发现 \textbf{双命名空间}:
\begin{itemize}[leftmargin=*]
  \item 层级路径: /CY1C7K/G138VS405/... (油田泵油单元, 数千条, 缓存到文件)
  \item 数字 ID: 5000, 5500, ... (保护装置, 14 块, 动态解析)
\end{itemize}

\section{最终架构全貌 (2026-07-13 定稿)}

\subsection{36 IOMan 实例连接目标 (netstat 实锤)}
\begin{table}[H]\centering
\begin{tabular}{rlll}
\toprule
\textbf{PID} & \textbf{连接目标} & \textbf{协议} & \textbf{说明} \\
\midrule
4472 & 11.66.12.130:8889 & pSpace psAPI & 唯一走 pSpace 的实例 \\
13852 & 172.23.9.3 & OPC DA DCOM & Kepware \#1 \\
8320 & 172.23.9.23 & OPC DA DCOM & Kepware \#2 \\
13776 & 172.26.6.3 & OPC DA DCOM & OPC Server \#3 \\
18044 & 172.23.18.194 & OPC DA DCOM & OPC Server \#4 \\
7616 & 172.21.14.192 & OPC DA DCOM & OPC Server \#5 (新发现) \\
4288 & CommBridge & Modbus RTU & 串口透传 \\
其余 29 个 & (无 TCP) & -- & 空闲/待命 \\
\bottomrule
\end{tabular}
\end{table}

\subsection{三系统数据架构}
\begin{verbatim}
系统A: 油田泵油 (层级路径, 30 TagID文件)
  CY1C7K: 134台设备·1922 Tag (G138VS405/D72VS304...)
  CY1C8K: ~170台设备 (132KB)
  pPluse: 数千Tag (脉冲量)
  访问: psAPI Tag Query + Oracle管线 (dgiot已同步26万条)

系统B: 保护装置 (数字ID, 14块)
  ID 5000-11000: IOMan#7 (IO7CEEB98) → pSpace
  设备: 02204060029-02204060123 (9台确认)
  访问: psAPI ReadList (自主采集已通)

系统C: OPC DA DCOM (5台Kepware, 50设备)
  172.23.9.3/9.23, 172.26.6.3, 172.23.18.194, 172.21.14.192
  访问: DCOM Browse (WinRM权限受限, 待解决)
\end{verbatim}

\subsection{TagID 文件格式 (逆向确认)}
\begin{verbatim}
[4B: uint32 count]
Repeated:
  [4B: uint32 strlen]
  [strlen bytes: ASCII path]
  [4B: uint32 index]
例: /CY1C7K/G138VS405/JD010707G138VS405ADL
     ─┬──  ───┬───  ─────────┬──────────
     站      设备          Tag名(ADL=相电流A...)
\end{verbatim}

\subsection{OPC DA 全通验证 (2026-07-13 22:11)}

\textbf{完整链路}:
\begin{verbatim}
Oracle PROJECT_TAGPAR (46,046 Tag)
  → LONGNAME: /CY1C8K/B1VD4VF20/JD010810B1VD4VF20ADL
    ↓ psAPI_Tag_GetIdByLongName
  → Tag ID: 5657
    ↓ psAPI_Real_ReadList
  → Value: 2.2017 @ 22:01:10 GOOD ✅
\end{verbatim}

\textbf{实测数据}:
\begin{table}[H]\centering
\begin{tabular}{rlll}
\toprule
\textbf{Tag ID} & \textbf{值} & \textbf{时间戳} & \textbf{来源} \\
\midrule
5001 & 1.8750 & 22:09:37 & 保护装置 (已知) \\
5657 & 2.2017 & 22:01:10 & OPC\_FC\_Client (新验证!) \\
11945 & 0.0000 & OLD & DX1ZRZ注水站 (离线) \\
\bottomrule
\end{tabular}
\end{table}

\textbf{结论}: 全部 292 设备 · 46,046 Tag 均可通过此链路自主读取实时数据。
OPC DA 完全攻克。

\section{DLAS 知识本体}

按 Data-Logic-Action-Security 四层框架组织全部逆向知识:

\begin{table}[H]\centering
\begin{tabular}{lr}
\toprule
\textbf{层} & \textbf{实体/关系/约束数} \\
\midrule
Data 层 & 60 设备·25 测点·14 通道·6 数据源 \\
Logic 层 & 100 关系·26 约束·12 设备类型 \\
Action 层 & 5 OPC 路径·ReadList 接口·Tag 扫描 \\
Security 层 & 凭据管理·生产网规则·审计建议 \\
\bottomrule
\end{tabular}
\end{table}

交付物: \texttt{实体清单.xlsx} (38 实体)、\texttt{关系矩阵.xlsx} (14 关系)、\texttt{约束规则库.xlsx} (18 规则)、\texttt{parse.db} (运行时本体)。

\section{安全审计}

\begin{table}[H]\centering
\begin{tabular}{lll}
\toprule
\textbf{发现} & \textbf{风险} & \textbf{建议} \\
\midrule
psAPISDK 调试日志泄漏 & 信息泄漏 & 生产环境禁用 DEBUG \\
IOMan.exe 明文密码 & 凭据泄漏 & 加密存储 \\
pSpace 无 TLS & 中间人攻击 & 启用 psAPI\_Server\_ConnectSSL \\
共享内存可读 & 数据泄漏 & 限制访问权限 \\
DCOM 135 端口暴露 & RCE & 防火墙限制 \\
\bottomrule
\end{tabular}
\end{table}

\section{工具清单}

\begin{table}[H]\centering
\begin{tabular}{lp{8cm}}
\toprule
\textbf{工具} & \textbf{用途} \\
\midrule
\texttt{dev\_env.py} & 一键启动全栈模拟 (7 服务 + N 设备) \\
\texttt{pspace\_collector.py} & pSpace 实时数据采集器 (32bit Python) \\
\texttt{scan\_tags.py} & Tag ID 批量扫描 (子进程隔离 + 限流) \\
\texttt{mock\_opc\_server.py} & Modbus TCP 模拟服务器 \\
\texttt{simulate\_131\_fullscale.py} & 全量设备 1:1 模拟 (56 设备) \\
\texttt{run\_commbridge\_local.py} & CommBridge 全链路测试 \\
\texttt{loop\_131\_monitor.py} & 131 生产服务器巡检 \\
\texttt{loop\_dgiot\_platform.py} & 本地平台健康巡检 \\
\bottomrule
\end{tabular}
\end{table}

\section{文档体系}

\begin{table}[H]\centering
\begin{tabular}{llc}
\toprule
\textbf{编号} & \textbf{文档} & \textbf{页} \\
\midrule
A & IO 服务器自主 OPC DA 采集技术方案 & 7 \\
B & pSpace 协议逆向研究报告 (v1.1 含 Tag 图谱) & 7 \\
C & 离线开发环境用户手册 & 7 \\
D & IO 服务器 DLAS 本体 v4.0 & 6 \\
E & OPC DA 点位自动发现方法论 & 6 \\
G & 本报告——逆向工程最终报告 & 9 \\
\bottomrule
\end{tabular}
\end{table}

\section{结论}

本项目通过三周逆向工程，完整还原了力控 pSpace 6.0 系统的 IO 采集架构。
核心突破在于：\textbf{自主发现 pSpace 连接凭证和 API 调用方式}，
实现了不依赖原厂 IoProject/IoMonitor 组件的自主 OPC DA 数据读取。

\textbf{建议后续工作}：
\begin{enumerate}[leftmargin=*]
  \item OPC Browse 方式一次性获取 Kepware 全量 tag tree
  \item 完成 8 个 IOMan 全部 80 设备的 Tag ID 映射
  \item psAPI\_Real\_NewSubscribeAndRead 回调机制集成
  \item 本机 pSpace mock 与真实 IoMonitor.exe 对接验证
  \item 生产网权限申请后启动自定义 IOMan 注入验证
\end{enumerate}

\end{document}
